En una zona de l'espai hi ha situades dues càrregues elèctriques puntuals de $3,0\ \mathrm{\mu C}$ i $-7,0\ \mathrm{\mu C}$ separades 15~cm l'una de l'altra. Calculeu:

\begin{apartat}{1}
El camp elèctric en un punt situat sobre la línia que uneix les càrregues. Aquest punt està situat a una distància de 5,0~cm de la càrrega de $3,0\ \mathrm{\mu C}$ i a 10~cm de la càrrega de $-7,0\ \mathrm{\mu C}$. Indiqueu-ne el mòdul, la direcció i el sentit.
\end{apartat}

\begin{apartat}{1}
El punt entre les dues càrregues en el qual el potencial és nul.
\end{apartat}

\blocetiquetat{Dada:}{%
  $k = \dfrac{1}{4\pi\varepsilon_0} = 8,99\times 10^{9}\ \mathrm{N\,m^2\,C^{-2}}$.}
